\subsection{A big divisor calibrated by the degree on the base}
\label{subsec:rank:calibration}

To compare orders on fibres with slopes on the base, we need a linear
functional on divisor classes of $X$ that agrees with base degree on
pullbacks and is nonnegative on effective divisors.  We construct this
functional together with a big divisor whose volume root and functional
value are controlled by one parameter.  The base polarization alone has
zero volume on $X$, whereas an arbitrary ample divisor need not satisfy
the required degree bound.

\begin{lemma}[Calibration]\label{lem:rank:calibration}
Let $h:X\to Y$ be a contraction of normal projective complex varieties,
where $X$ is a $\Q$-factorial canonical $n$-fold with $K_X\sim0$ and
irregularity zero on a resolution.  Suppose that $0<b:=\dim Y<n$ and
that no prime divisor of $X$ maps into a subset of codimension at least
two in $Y$.  Let $\widehat Y$ be a normal projective variety isomorphic
to $Y$ in codimension one, and let $H$ be a very ample divisor on
$\widehat Y$ satisfying
\[
 H^b\le N_0,
 \qquad
 \bigl|K_{\widehat Y}\cdot H^{b-1}\bigr|\le N_0.
\]
Pullback from the common smooth open subsets of the bases defines
\[
 W=h^*\Cl(Y)_{\R}
   =h^*\Cl(\widehat Y)_{\R}
   \subset N^1(X)_{\R},
 \qquad
 \ell(h^*L)=L\cdot H^{b-1}.
\]
Let $A=h^*H$, represented by an effective integral Weil divisor.
There is a linear extension
$\ell_X:N^1(X)_{\R}\to\R$ of $\ell$, nonnegative on the
pseudo-effective cone, and there are a rational number $\delta>0$ and
a big rational divisor $D$ such that
\begin{equation}\label{eq:rank:calibration}
 c\delta\le v(D)\le C\delta,
 \qquad
 D-\frac{\delta}{3}A\text{ is big},
 \qquad
 \ell_X(D)\le C\delta.
\end{equation}
Here $v(D)=\vol(D)^{1/n}$, $c>0$ depends only on $n$, and $C$ depends
only on $n,N_0$.
\end{lemma}

The two conditions beyond $v(D)\asymp\delta$ have distinct roles.
The bigness of $D-\delta A/3$ will force the low-order families to be
vertical.  The bound on $\ell_X(D)$ will convert cotangent
semipositivity into a slope bound on the base.

\begin{proof}
We first construct a supporting functional for the pseudo-effective
cone and then choose a point on which volume has the required scale.
The hypothesis on irregularity, together with $\Q$-factoriality, gives
\[
 \Cl(X)_{\R}=\Pic(X)_{\R}=N^1(X)_{\R}.
\]
Indeed pullback to a resolution detects numerically trivial line
bundles, whose classes are torsion when the Picard variety is zero.
We can therefore pass between real linear equivalence and numerical
classes throughout the proof.

Choose common smooth big open subsets of $Y$ and $\widehat Y$.
Their inverse image in $X$ has complement of codimension at least two,
so Cartier pullbacks there extend uniquely as Weil divisors on $X$.
This defines the indicated pullback on class groups.  It is injective:
if $s=n-b$ and $Q$ is an ample Cartier divisor on $X$, the operation
\[
 B\longmapsto h_*(B\cdot Q^s)
\]
on divisor classes is a left inverse to pullback, up to multiplication
by the positive degree of $Q^s$ on the generic fibre.  The same
operation sends effective classes to effective classes or zero, by
using general representatives of sufficiently divisible multiples of
$Q$.

The degree $\ell$ is strictly positive on
$W\cap\overline{\operatorname{Eff}}(X)\setminus\{0\}$.
To see the point requiring care, choose a finite linear projection
$f:\widehat Y\to\PP^b$ with $f^*\OO_{\PP^b}(1)=\OO_{\widehat Y}(H)$.
For every effective real Weil divisor $L$ on $\widehat Y$,
\begin{equation}\label{eq:rank:degree-dominates}
 \ell(L)H-L\quad\text{has an effective representative}.
\end{equation}
In fact $f^*f_*L-L\ge0$ and
$f_*L\sim_{\R}\ell(L)\OO_{\PP^b}(1)$.
The intersect-and-push operation just described shows that if a
pullback class is pseudo-effective on $X$, its base class lies in the
closure of effective base classes.  Applying
\eqref{eq:rank:degree-dominates} to approximating classes and pulling
back shows that a degree-zero such class and its negative are both
pseudo-effective on $X$.  The pseudo-effective cone in $N^1(X)_{\R}$
contains no line, so the class is zero.  This also proves that the
formula defining $\ell$ is well defined on $W$.

Fix an ample class $P_0$ on $X$, and minimize $\ell(L)$ on the set
\[
 \{L\in W:P_0+L\in\overline{\operatorname{Eff}}(X)\}.
\]
The set is nonempty.  Its subsets on which $\ell$ is bounded above
are bounded: otherwise, dividing an unbounded sequence by its norm
and passing to a limit would give a nonzero pseudo-effective class in
$W$ of nonpositive degree.  Closedness therefore gives a minimum
$c_0=\ell(L_0)$.  Put $d_0=P_0+L_0$.
The open big cone is disjoint from
\[
 P_0+\{L\in W:\ell(L)\le c_0\}.
\]
Indeed a big point in this affine half-space would remain big after
subtracting a sufficiently small positive multiple of $A$, contradicting
minimality.  Separating these two convex sets gives a nonzero linear
functional $\lambda$, nonnegative on the pseudo-effective cone, with
$\lambda(d_0)=0$.  Its restriction to $W$ is a nonnegative multiple
of $\ell$: it vanishes on $\ker\ell$, and the orientation follows
from the half-space being separated.  That multiple is positive,
since otherwise $\lambda(P_0)=\lambda(d_0)=0$, whereas a nonzero
nonnegative functional is strictly positive at an interior point of
the cone.  Normalize $\lambda$ to obtain $\ell_X$.

For $t>0$, set $d_t=d_0+tA$.  These classes are big.  For large $t$,
$L_0+tA$ is pseudo-effective, since sufficiently large multiples of
$H$ added to any fixed real Weil class on the base are effective;
convexity with the pseudo-effective class $d_0$ then proves the
assertion for every $t>0$.  On the other hand, $A$ is not big: its
restriction to the positive-dimensional generic fibre is trivial.
Continuity and homogeneity of volume give
\begin{equation}\label{eq:rank:large-calibration-parameter}
 \frac{v(d_t)}t=v(A+t^{-1}d_0)\longrightarrow0
 \qquad(t\longrightarrow\infty),
 \qquad
 \ell_X(d_t)=tH^b.
\end{equation}
It remains to bound this volume ratio from below for small $t$.
This is where the trivial canonical class enters the construction.

Let $p:\widetilde X\to X$ be a smooth projective resolution and write
the divisorial Zariski decomposition of
$\widetilde d_0=p^*d_0$ as
\[
 \widetilde d_0=P_\sigma(\widetilde d_0)
                      +N_\sigma(\widetilde d_0).
\]
We use its numerical formulation, as in
\citep[Chapter~III]{Nakayama2004} and \citep{Boucksom2004}.
Fix the dimension-dependent integer $M=M(n)$ in Birkar's
Calabi--Yau effective birationality theorem
\citep[Corollary~1.4]{Birkar2023}.
Suppose first that $P_\sigma(\widetilde d_0)$ is numerically nonzero.
For an ample class $J$ on $\widetilde X$, monotonicity of positive
intersection products gives
\begin{align*}
 \vol(\widetilde d_0+uJ)
 &=\bigl\langle(\widetilde d_0+uJ)^n\bigr\rangle\\
 &\ge u^{n-1}
 P_\sigma(\widetilde d_0+uJ)\cdot J^{n-1}.
\end{align*}
Here the trace of the one-class positive product is the positive
part of the divisorial Zariski decomposition; see
\citep[Definition~2.10, Proposition~2.13, Theorem~3.1, and
Section~3.4]{BFJ2009}.
As $u\downarrow0$, the intersection on the right tends to
$P_\sigma(\widetilde d_0)\cdot J^{n-1}>0$.
Choose $a>0$ such that $p^*d_1-aJ$ is pseudo-effective.  For
$0<t<1$, the identity
$d_t=(1-t)d_0+td_1$ and volume monotonicity yield
\[
 \vol(d_t)\ge(1-t)^n
 \vol\left(\widetilde d_0+\frac{at}{1-t}J\right).
\]
Consequently $v(d_t)/t\to\infty$ as $t\downarrow0$.
The constants in this divergence may depend on $X$; no uniform
threshold in $t$ is needed.

Suppose instead that $P_\sigma(\widetilde d_0)$ is numerically zero.
Then $d_0$ has a unique effective real representative $N_0'$.
For existence, push forward $N_\sigma(\widetilde d_0)$; numerical
and real linear equivalence agree here.  For uniqueness, the pullback
of any effective representative of $d_0$ dominates
$N_\sigma(\widetilde d_0)$.  Their difference is effective and
numerically zero, hence zero.  Let $E_0$ be the reduced support of
$N_0'$.  It has Kodaira dimension zero: if some multiple of $E_0$ had
a different effective representative, replacing a sufficiently small
positive multiple of $E_0$ inside $N_0'$ would contradict uniqueness.

Choose a small positive rational $\varepsilon$ for which
$(X,\varepsilon E_0)$ is klt.  Theorem~\ref{input:good-model}
gives a terminating MMP for this pair and a
good model $\overline X$.  Precisely the components of $E_0$ are
contracted.  Any divisorial step must contract a component of its
transform, since $K\sim0$ and a curve outside the effective boundary
has nonnegative intersection with it.  Conversely, semiampleness on
the good model and $\kappa(E_0)=0$ force the effective pushforward
of $E_0$ to vanish.  These steps are crepant for the zero boundary, so
$\overline X$ is again $\Q$-factorial canonical with principal
canonical divisor.  The pushforward $\overline A$ of $A$ is an
integral Weil divisor and is big: the big class $d_t$ pushes forward
to $t\overline A$.  Birkar's theorem now gives
\begin{equation}\label{eq:rank:integral-volume-lower-bound}
 \vol(\overline A)\ge M^{-n}
 \qquad\text{\citep[Corollary~1.4]{Birkar2023}}.
\end{equation}
Its hypotheses are exactly that $\overline X$ is projective klt,
$K_{\overline X}\sim_{\R}0$, and $\overline A$ is big and integral;
no Cartier-index bound is required.

A single rational number $u\ge0$ gives inclusions
\[
 H^0(\overline X,m\overline A)
 \subseteq H^0(X,m(A+uE_0))
\]
for every sufficiently divisible $m$.  Indeed, on a common resolution
pull back $\overline A$ and push forward to $X$.  The resulting
rational divisor differs from $A$ only on $E_0$, and its difference
is bounded above by $uE_0$.  Thus
$\vol(A+uE_0)\ge\vol(\overline A)$.
For all sufficiently small $t>0$ we have $N_0'\ge tuE_0$, and hence
\[
 \vol(d_t)=\vol(N_0'+tA)
 \ge t^n\vol(A+uE_0)\ge t^n M^{-n}.
\]
The same conclusion includes $E_0=0$, with $u=0$.

In both cases the ratio $v(d_t)/t$ is at least $M^{-1}$ for all
sufficiently small $t>0$.  By
\eqref{eq:rank:large-calibration-parameter} and continuity, it takes
the value $(2M)^{-1}$ at some $t>0$.  At this parameter,
\[
 d_t-\frac t3A=d_{2t/3}\text{ is big},
 \qquad \ell_X(d_t)=tH^b\le N_0t.
\]
Approximate $t$ by a positive rational number $\delta$ and $d_t$ by a
rational class $D$.  Openness of the big cone and continuity of volume
and $\ell_X$ preserve these assertions with, for example, lower
constant $1/(4M)$ and suitable upper constant depending only on
$n,N_0$.  Since $X$ is $\Q$-factorial, the rational class has a
rational Cartier divisor representative.  This proves
\eqref{eq:rank:calibration}.
\end{proof}
